\documentclass[11pt]{article}
\usepackage[a4paper,margin=1in]{geometry}
\usepackage{amsmath,amssymb,amsthm}
\usepackage{graphicx}
\usepackage{microtype}
\usepackage{xcolor}
\usepackage[hidelinks]{hyperref}
\usepackage{enumitem}

\newtheorem{theorem}{Theorem}
\newtheorem{lemma}[theorem]{Lemma}
\newtheorem{proposition}[theorem]{Proposition}
\newcommand{\Space}{\operatorname{Space}}
\newcommand{\Opp}{\operatorname{Op}^{<\omega}}
\newcommand{\Id}{\operatorname{Id}}
\newcommand{\ran}{\operatorname{ran}}
\newcommand{\cH}{\mathcal H}
\newcommand{\cJ}{\mathcal J}

\title{No largest operator ideal with Hilbert identity spaces\\
\large A candidate full negative answer to Question 6.2 of arXiv:2005.07672}
\author{Automated mathematics run \texttt{fa\_banach\_001}, lane 5}
\date{9 August 2026}

\begin{document}
\maketitle

\begin{center}
\fbox{\parbox{0.91\textwidth}{\textbf{Status.}
Candidate full solution, likely valid, requiring expert human review.
The argument answers the question negatively over both real and complex
scalars.}}
\end{center}

\section{Source question}

Ferenczi's paper \cite{Ferenczi} asks on page 22, Question 6.2, whether
there is a largest operator ideal whose identity spaces are exactly the
separable Hilbert spaces.  Here is the source passage.

\begin{figure}[h]
  \centering
  \includegraphics[width=\textwidth]{figures/open_problem_crop.png}
  \caption{Question 6.2 and its context in \cite[p.~22]{Ferenczi}.}
\end{figure}

Writing $\cH$ for the space ideal of separable Banach spaces isomorphic to a
Hilbert space, the question is:
\begin{quote}
Does there exist a largest operator ideal $\mathcal U$ with
$\cH=\Space(\mathcal U)$?
\end{quote}

\section{Definitions and imported facts}

Fix one scalar field $\mathbb K\in\{\mathbb R,\mathbb C\}$ and put
$H=\ell_2(\mathbb K)$. For a Banach space $E$, let $\Opp(E)$ be the operator
ideal consisting of all operators which factor through $E^n$ for some
$n\geq 1$. If $\mathcal U$ is an operator ideal, then
\[
  \Space(\mathcal U)=\{Z:\Id_Z\in\mathcal U\}.
\]

We use four established inputs.
\begin{enumerate}[label=(\roman*)]
\item Gowers and Maurey's shift space $X=X_S(\mathbb K)$ contains no
unconditional basic sequence \cite[Section~5.2]{GM}. In particular, it
contains no infinite-dimensional Hilbert space.
\item Ferenczi constructs, for $t\in\{1,-1\}$, a compact operator $K_t$ such
that
\[
 T_t=\Id_X-tS+K_t,\qquad Y_t=\ran T_t,
\]
where $Y_t$ has infinite codimension in $X$ \cite[Section~4]{Ferenczi}.
Moreover, if $m,n\geq1$, then no infinite-dimensional complemented subspace
of $X^m$ is isomorphic to a subspace of $Y_t^n$ \cite[Proposition~4.4]{Ferenczi}.
\item If $E$ and $F$ are essentially incomparable, then every complemented
subspace of $E\oplus F$ is isomorphic to $E_0\oplus F_0$ for complemented
subspaces $E_0\subseteq E$ and $F_0\subseteq F$ \cite{Gonzalez}.
\item Strictly singular perturbations preserve the upper semi-Fredholm
class. Equivalently, if $A$ is upper semi-Fredholm and $R$ is strictly
singular, then $A+R$ is upper semi-Fredholm \cite[Chapter~2]{LT}.
\end{enumerate}

Recall also that strictly singular operators are inessential. Thus, if every
operator in both directions between two spaces is strictly singular, those
spaces are essentially incomparable.

\section{Proof intuition}

Ferenczi destroys a hypothetical largest \emph{proper} ideal using two
factorization ideals associated with $X$ and with $Y_t$. To preserve exactly
the Hilbert identity spaces, adjoin a common Hilbert summand before repeating
the construction:
\[
 \cJ_t=\Opp(H\oplus X)\cap\Opp(H\oplus Y_t).
\]
The common $H$ makes every Hilbert identity belong to $\cJ_t$. The point is
that it adds no other identity. Indeed, a space whose identity belongs to
$\cJ_t$ is complemented in both $H\oplus X^m$ and $H\oplus Y_t^n$.
Essential incomparability separates each complemented copy into a Hilbert
piece and an exotic piece. Comparing the two decompositions forces the exotic
piece, modulo a finite-dimensional kernel, to embed from $X^m$ into $Y_t^n$.
Ferenczi's rigidity theorem rules this out. Hence the common identity spaces
are exactly Hilbertian, while the two operators $T_1,T_{-1}$ still belong to
the corresponding ideals and still have a Fredholm sum.

\section{The common-complemented-subspace argument}

\begin{lemma}\label{lem:nohilbert}
For every $m,n\geq1$, neither $X^m$ nor $Y_t^n$ contains a subspace
isomorphic to $H$.
\end{lemma}

\begin{proof}
The space $X$ contains no unconditional basic sequence, hence no copy of
$H$. Suppose $H$ embedded into $X^m$. Each coordinate of the embedding is
strictly singular: otherwise it would be bounded below on an
infinite-dimensional subspace of $H$, and every such subspace is isomorphic
to $H$, producing a copy of $H$ in $X$. The column embedding, being a finite
sum of strictly singular operators after the coordinate inclusions, would
then be strictly singular, contradicting that it is bounded below. Thus
$X^m$ contains no $H$. Since $Y_t^n$ is a subspace of $X^n$, the second
assertion follows.
\end{proof}

\begin{lemma}\label{lem:decompose}
For each $m,n\geq1$, the pairs $(H,X^m)$ and $(H,Y_t^n)$ are essentially
incomparable. Consequently every complemented subspace of $H\oplus X^m$ is
isomorphic to $H_0\oplus E_0$, where $H_0$ is Hilbertian and $E_0$ is
complemented in $X^m$; the analogous assertion holds for $H\oplus Y_t^n$.
\end{lemma}

\begin{proof}
If an operator $A:H\to X^m$ were not strictly singular, it would be bounded
below on an infinite-dimensional subspace of $H$, hence would put a copy of
$H$ in $X^m$, contrary to Lemma~\ref{lem:nohilbert}. If
$B:X^m\to H$ were not strictly singular, the subspace on which it is bounded
below would itself be Hilbertian, giving the same contradiction. Thus all
cross operators are strictly singular and therefore inessential. The same
argument applies to $Y_t^n$. The stated decompositions now follow from
Gonzalez's theorem \cite{Gonzalez}.
\end{proof}

\begin{lemma}\label{lem:common}
If a Banach space $Z$ is isomorphic to a complemented subspace of both
$H\oplus X^m$ and $H\oplus Y_t^n$, then $Z$ is Hilbertian.
\end{lemma}

\begin{proof}
By Lemma~\ref{lem:decompose}, there are decompositions
\[
 Z\simeq H_0\oplus E_0\simeq H_1\oplus F_0,
\]
where $H_0,H_1$ are Hilbertian, $E_0$ is complemented in $X^m$, and $F_0$
is complemented in $Y_t^n$. Let
$R:H_0\oplus E_0\to H_1\oplus F_0$ be an isomorphism. On $E_0$, write
\[
 R|_{E_0}=(A,B),\qquad A:E_0\to H_1,\quad B:E_0\to F_0.
\]
The operator $A$ is strictly singular. Indeed, otherwise $A$ would be
bounded below on an infinite-dimensional subspace of $E_0$, making that
subspace Hilbertian and contradicting Lemma~\ref{lem:nohilbert}.

The column $(A,B):E_0\to H_1\oplus F_0$ is bounded below, hence upper
semi-Fredholm. It is a strictly singular perturbation of $(0,B)$, so $(0,B)$,
and equivalently $B$, is upper semi-Fredholm. Therefore $B$ has finite
dimensional kernel and closed range. Choose a closed complement $E_1$ of
$\ker B$ in $E_0$. Then $E_1$ is finite codimensional in $E_0$, is
complemented in $X^m$, and $B|_{E_1}$ is an isomorphic embedding into
$F_0\subseteq Y_t^n$. Ferenczi's Proposition~4.4 forces $E_1$ to be finite
dimensional. Hence $E_0$ is finite dimensional, and
$Z\simeq H_0\oplus E_0$ is Hilbertian.
\end{proof}

\section{Full negative answer}

\begin{theorem}\label{thm:main}
Over either the real or complex scalars, there is no largest operator ideal
$\mathcal U$ satisfying $\Space(\mathcal U)=\cH$.
\end{theorem}

\begin{proof}
For $t\in\{1,-1\}$ define
\[
 \cJ_t:=\Opp(H\oplus X)\cap\Opp(H\oplus Y_t).
\]
We first show
\begin{equation}\label{eq:space}
 \Space(\cJ_t)=\cH.
\end{equation}
Every separable Hilbert space is finite dimensional or isomorphic to $H$, so
its identity factors through the common $H$ summand and belongs to $\cJ_t$.
Conversely, suppose $\Id_Z\in\cJ_t$. A factorization of an identity through a
space realizes its domain as a complemented subspace of that space. Hence,
for some $m,n\geq1$, the space $Z$ is isomorphic to a complemented subspace of
both
\[
 (H\oplus X)^m\simeq H\oplus X^m
 \quad\text{and}\quad
 (H\oplus Y_t)^n\simeq H\oplus Y_t^n.
\]
Lemma~\ref{lem:common} says that $Z$ is Hilbertian. This proves
\eqref{eq:space}.

Assume now that a largest ideal $\mathcal U$ with
$\Space(\mathcal U)=\cH$ exists. By \eqref{eq:space}, it contains both
$\cJ_1$ and $\cJ_{-1}$. The operator $T_t:X\to Y_t$, also viewed as an
operator $X\to X$, factors through $X$ and through $Y_t$. It therefore
factors through both $H\oplus X$ and $H\oplus Y_t$, so $T_t\in\cJ_t$.
Consequently $T_1,T_{-1}\in\mathcal U$ and
\[
 F:=T_1+T_{-1}=2\Id_X+K_1+K_{-1}\in\mathcal U(X,X).
\]
The operator $F$ is Fredholm. Choose a parametrix $R$ with
$RF=\Id_X-Q$ for a finite-rank operator $Q$. Ideality gives $RF\in\mathcal U$,
and every operator ideal contains $Q$; hence $\Id_X=RF+Q\in\mathcal U$.
Thus $X\in\Space(\mathcal U)=\cH$, impossible because $X$ contains no
unconditional basic sequence and so is not Hilbertian. This contradiction
proves the theorem.
\end{proof}

\section{Verification and limitations}

The proof is non-computational. The exact source question, the construction
of $Y_t$, Proposition 4.4, and the Fredholm identity for $T_1+T_{-1}$ were
checked against the source TeX and PDF. The no-unconditional-basic-sequence
statement was checked in the primary Gowers--Maurey source
\cite[Section~5.2]{GM}, and the direct-sum theorem against Gonzalez's article
record and abstract \cite{Gonzalez}.

A bounded novelty search on 9 August 2026 covered the run registry, cheap
arXiv source indexes, and indexed web/arXiv results for the exact question and
the phrases ``largest operator ideal'', ``Hilbert space ideal'', and
``$\Space(\mathcal U)=\cH$''. It found the source paper and later citations of
its proper-ideal theorem, but no claimed resolution of Question 6.2. This is
not exhaustive, so novelty confidence is moderate. The result should be cited
only as a candidate full solution until an expert independently checks the
three imported structural facts, especially the semi-Fredholm perturbation
step in Lemma~\ref{lem:common}.

\begin{thebibliography}{9}

\bibitem{Ferenczi}
V. Ferenczi,
\emph{There is no largest proper operator ideal},
arXiv:2005.07672v2; Math. Ann. 387 (2023), 1043--1072.

\bibitem{GM}
W. T. Gowers and B. Maurey,
\emph{Banach spaces with small spaces of operators},
Math. Ann. 307 (1997), 543--568; arXiv:math/9407209.

\bibitem{Gonzalez}
M. Gonzalez,
\emph{On essentially incomparable Banach spaces},
Math. Z. 215 (1994), 621--629.

\bibitem{LT}
J. Lindenstrauss and L. Tzafriri,
\emph{Classical Banach Spaces I and II},
Springer, 1996 reprint.

\end{thebibliography}

\end{document}
